\section{Практичне дослідження узгодження спільного секрету}

\subsection{Вибір криптографічного механізму та бібліотеки}
Для проведення експерименту з узгодження спільного секрету (Key Agreement) обрано механізм X25519,
який базується на еліптичній кривій Curve25519. Цей протокол є сучасним стандартом (RFC 7748), оскільки
забезпечує високу швидкодію, стійкість до атак по побічних каналах та високий рівень безпеки (128 біт). 

Для реалізації використано бібліотеку \texttt{cryptography} мови Python, зокрема модуль
\texttt{cryptography.hazmat.primitives.asymmetric.x25519}.

\subsection{Генерація ключів Alice}
Аліса генерує власну пару ключів X25519 (приватний та відповідний йому відкритий ключ).

\begin{lstlisting}[language=Python]
from cryptography.hazmat.primitives.asymmetric import x25519
from cryptography.hazmat.primitives import serialization

alice_private_key = x25519.X25519PrivateKey.generate()
alice_public_key = alice_private_key.public_key()

print("Alice key pair generated")
\end{lstlisting}

Результат виконання:
\begin{lstlisting}
Alice key pair generated
\end{lstlisting}

\subsection{Генерація ключів Bob}
Боб незалежно генерує власну пару ключів X25519.

\begin{lstlisting}[language=Python]
bob_private_key = x25519.X25519PrivateKey.generate()
bob_public_key = bob_private_key.public_key()

print("Bob key pair generated")
\end{lstlisting}

Результат виконання:
\begin{lstlisting}
Bob key pair generated
\end{lstlisting}

\subsection{Обмін відкритими ключами}
Учасники експортують свої відкриті ключі у байтовий формат для подальшої передачі через недовірену мережу.

\begin{lstlisting}[language=Python]
alice_pub_bytes = alice_public_key.public_bytes(
    encoding=serialization.Encoding.Raw,
    format=serialization.PublicFormat.Raw
)
bob_pub_bytes = bob_public_key.public_bytes(
    encoding=serialization.Encoding.Raw,
    format=serialization.PublicFormat.Raw
)

print("\nPublic keys exchanged via network:")
print("Alice Public Key (hex):", alice_pub_bytes.hex())
print("Bob Public Key (hex):  ", bob_pub_bytes.hex())
\end{lstlisting}

Результат виконання:
\begin{lstlisting}
Public keys exchanged via network:
Alice Public Key (hex): 8b1a30d12461b22fa4b51740dd7edee7ee9fbfa205896c4acfce331e551e3775
Bob Public Key (hex):   d56e6bc0a7c316ad64557ccdfc04d86bdbac162c8e427972c47d0561be703805
\end{lstlisting}

\subsection{Формування спільного секрету Alice}
Аліса обчислює спільний секрет, використовуючи власний приватний ключ та відкритий ключ Боба.

\begin{lstlisting}[language=Python]
alice_shared_secret = alice_private_key.exchange(bob_public_key)

print("\nAlice derived shared secret:")
print("Alice Secret (hex):", alice_shared_secret.hex())
\end{lstlisting}

Результат виконання:
\begin{lstlisting}
Alice derived shared secret:
Alice Secret (hex): c75a40118b93810d7cfbcdb15267733444be22e913bd60e025c1146ae9b65665
\end{lstlisting}

\subsection{Формування спільного секрету Bob}
Боб обчислює спільний секрет, використовуючи власний приватний ключ та відкритий ключ Аліси.

\begin{lstlisting}[language=Python]
bob_shared_secret = bob_private_key.exchange(alice_public_key)

print("\nBob derived shared secret:")
print("Bob Secret (hex):  ", bob_shared_secret.hex())
\end{lstlisting}

Результат виконання:
\begin{lstlisting}
Bob derived shared secret:
Bob Secret (hex):   c75a40118b93810d7cfbcdb15267733444be22e913bd60e025c1146ae9b65665
\end{lstlisting}

\subsection{Порівняння отриманих спільних секретів}
Перевіряємо ідентичність обчислених секретів обома сторонами.

\begin{lstlisting}[language=Python]
is_equal = (alice_shared_secret == bob_shared_secret)

print("\nComparison result:")
print("Shared secrets are equal:", is_equal)
\end{lstlisting}

Результат виконання:
\begin{lstlisting}
Comparison result:
Shared secrets are equal: True
\end{lstlisting}

\subsection{Аналіз результатів}

\textbf{Чому Alice та Bob отримали однаковий спільний секрет?}

Аліса та Боб отримали однаковий секрет завдяки математичним властивостям операцій над точками
еліптичної кривої. Обчислення базується на комутативності скалярного множення:
$$K_{shared} = d_A \cdot (d_B \cdot G) = d_B \cdot (d_A \cdot G) = (d_A \cdot d_B) \cdot G$$
де $d_A$ та $d_B$ --- приватні ключі Аліси та Боба, а $G$ --- базова точка кривої X25519.

\textbf{Які дані передавалися між сторонами?}

Між сторонами через мережу передавалися виключно відкриті ключі (байтові послідовності
$A_{pub} = d_A \cdot G$ та $B_{pub} = d_B \cdot G$).

\textbf{Чи передавався сам спільний секрет через мережу?}

Ні, сам спільний секрет через мережу ніколи не передавався. Обидві сторони обчислили його повністю
незалежно та локально у власному середовищі.

\textbf{Чи може пасивний спостерігач сформувати цей секрет, маючи лише передані відкриті ключі?}

Ні, пасивний спостерігач не може сформувати секрет. Для цього йому потрібно знати хоча б один із
приватних ключів ($d_A$ або $d_B$). Знаходячись у мережі та бачачи лише $A_{pub}$ і $B_{pub}$, визначити
значення $d_A \cdot d_B \cdot G$ є обчислювально нездійсненною задачею, яка називається
обчислювальною задачею Діффі--Хеллмана (CDH) і спирається на складність дискретного логарифмування в
групі точок еліптичної кривої.

\textbf{Чому сам по собі механізм узгодження ключа не обов'язково підтверджує особу іншої сторони?}

Базовий протокол ECDH є анонімним і не забезпечує аутентифікацію сторін за замовчуванням. Відкритий
ключ X25519 сам по собі не містить інформації про особу власника. Це робить протокол вразливим до атак
типу <<людина посередині>> (Man-in-the-Middle), коли зловмисник може перехопити та підмінити відкриті
ключі під час обміну. Для підтвердження особи відкриті ключі повинні додатково підписуватися цифровим
підписом або підтверджуватися сертифікатами X.509 у рамках інфраструктури відкритих ключів
(PKI)~\cite{RFC7748,NIST80057}.
